\documentclass[conference]{IEEEtran}
\usepackage{cite,amsmath,amssymb,graphicx,textcomp,booktabs,hyperref,url}

\title{When Code Passes Once but Fails Twice:\\[2pt] Variability and Semantic Brittleness in AI Coding Agents}

\author{\IEEEauthorblockN{Sun Chengxi}
\IEEEauthorblockA{Independent Researcher\\
3852236135@qq.com}
}

\begin{document}

\maketitle

\begin{abstract}
AI coding agents are typically evaluated by single-run pass rates on benchmark tasks.
However, a single pass does not guarantee reliable performance under repeated evaluation.
We present a case study where an AI coding agent (Claude Code with deepseek-v4-flash) achieved 100\% pass rate across 18 tasks in single-run evaluation, but a distributed lock bug-fixing task (H2) succeeded in only 3 out of 5 repeated runs---despite all five implementations being semantically correct.
We propose a three-layer failure model that distinguishes (L1) semantic correctness, (L2) interface compatibility, and (L3) temporal stability, and introduce a variance index for quantifying run-to-run reliability.
Our findings suggest that single-run evaluations may systematically overestimate the stability of AI coding agents, particularly for concurrent and distributed systems tasks.
\end{abstract}

\section{Introduction}

AI coding agents have achieved high pass rates on established benchmarks such as HumanEval~\cite{chen2021evaluating} and SWE-bench~\cite{jimenez2023swebench}.
These evaluations share a common methodology: each task is executed once, and the pass/fail outcome is reported as a binary result.
This single-run paradigm implicitly assumes that AI code generation is deterministic---an assumption that the stochastic nature of large language models (LLMs) contradicts~\cite{brown2020language}.
The concern is magnified in the agentic coding paradigm, where multi-turn interactions compound generation uncertainty at each step: a decision made in an early turn constrains the solution space of later turns, causing small initial variations to amplify into qualitatively different outcomes.
Recent benchmarks for LLM-based agents, including AgentBench~\cite{liu2023agentbench}, have begun to address multi-turn capability, yet they still aggregate results into single-point metrics that obscure task-specific variability.
As AI coding agents transition from controlled evaluations to production deployment, the gap between single-run pass rates and real-world reliability becomes a pressing concern.
A task that can pass under standard evaluation may fail in ways invisible to existing pipelines yet consequential in production---particularly for concurrent and distributed systems, where execution non-determinism compounds generation non-determinism.
The research community currently lacks a systematic vocabulary for categorizing and measuring these compound failures, which is the gap this paper aims to address.

Our case study of Claude Code with deepseek-v4-flash on 18 software development tasks reveals a critical gap in this evaluation paradigm.
While the agent achieved 100\% pass rate in single-run evaluation, repeated runs of a distributed lock bug-fixing task (H2) succeeded in only 3 out of 5 attempts---a 60\% pass rate with variance 2.56.
This variability was not due to logical errors; all five implementations were semantically correct.
Instead, failures occurred due to \textit{interface incompatibility} (struct field naming mismatches causing compilation errors) and \textit{temporal instability} (race conditions in concurrent goroutine execution).

We make the following contributions:

\begin{enumerate}
\item A \textbf{three-layer failure model} distinguishing semantic correctness, interface compatibility, and temporal stability---dimensions of AI code quality that single-run pass/fail metrics conflate.
\item A \textbf{variance index} for quantifying run-to-run reliability, supported by empirical data from 30 repeated runs across six tasks.
\item A detailed \textbf{case study} of a distributed lock task (H2) demonstrating all three failure layers.
\end{enumerate}
Taken together, these contributions provide a foundation for moving beyond single-run pass/fail metrics toward a multi-dimensional understanding of AI coding agent reliability.

\section{Background and Motivation}

\subsection{Limitations of Single-Run Evaluation}

Current benchmarks for AI coding agents rely almost exclusively on single-run pass rates.
HumanEval~\cite{chen2021evaluating} reports pass@k from independent samples but treats each sample independently, not as repeated measurements of the same task.
SWE-bench~\cite{jimenez2023swebench} evaluates on real GitHub issues but does not systematically analyze run-to-run variance.
Additional benchmarks such as APPS~\cite{hendrycks2021measuring} and MBPP~\cite{austin2021program} follow the same single-sample methodology, testing standalone function synthesis rather than repository-level agentic workflows.
Even pass@k---a metric designed to capture multiple independent attempts---treats each sample as a fresh draw from the same distribution rather than a repeated measurement of the same task, and thus reflects population-level performance rather than task-level consistency.
These benchmarks were originally designed for code synthesis from natural language descriptions, not for the iterative debugging, context navigation, and integration workflows that characterize modern agentic usage.
The prevailing assumption is that a task either passes or fails deterministically.

This assumption is increasingly problematic.
LLMs are stochastic by design---the same prompt can produce different outputs across invocations due to sampling temperature, top-p selection, and GPU non-determinism.
Common decoding strategies (greedy, temperature $T > 0$, top-$k$, top-$p$) directly control output diversity, yet most benchmark evaluations report results under a single configuration without analyzing its impact on downstream task performance.
In the specific case of code generation, small syntactic variations---variable naming, function signature ordering, import style---can cascade into semantically meaningful differences in multi-turn settings where the output of one turn becomes the prompt context for the next.
This phenomenon is fundamentally distinct from variability in natural language: in code, a minor deviation can produce a compilation error or a runtime failure, whereas in prose it would merely change stylistic register.
While prior work has examined output variability in natural language tasks, systematic analysis of variability in code generation for agentic systems remains limited.
The multi-turn nature of modern coding agents means that early-generation choices---such as which library to import or what function signature to adopt---propagate through subsequent steps, making variability both more likely and more consequential than in single-turn synthesis.
Existing benchmarks also share structural limitations in task design: they test isolated function synthesis or patch generation rather than the iterative debugging, test-driven development, and integration workflows that constitute real-world agent usage.
The growing adoption of AI coding assistants---from automated bug fixing to feature implementation to code review---demands evaluation frameworks that measure not only functional correctness but also behavioral consistency across repeated interactions.
Without such frameworks, the community risks deploying agents that perform well in controlled evaluations yet exhibit unpredictable behavior under production conditions.

\subsection{The Need for Multi-Dimensional Evaluation}

A coding agent produces not just a binary outcome, but an artifact with multiple quality dimensions:
\begin{itemize}
\item Does the code correctly implement the algorithm? (semantic)
\item Does the code integrate with the existing codebase? (interface)
\item Can the code consistently pass tests across multiple runs? (temporal)
\end{itemize}
Single-run pass rates collapse these dimensions into one bit of information, obscuring important failure modes unique to AI-generated code.

\section{Three-Layer Failure Model}

We propose that AI-generated code failures fall into three independent categories, forming a hierarchy of severity:

\subsection{L1: Semantic Correctness}

The algorithm correctly implements the required functionality and all test assertions pass.
This is what traditional pass/fail metrics measure.
An L1 failure means the code is simply wrong---the most severe and easily detected category.
Consider an agent tasked with implementing a priority queue: it may generate a correct linear-scan-based approach that passes basic test cases but fails under performance constraints when an O(n log n) heap-based solution is required.
This subclass of L1 failure---algorithmic adequacy masked by insufficient test coverage---is particularly misleading because standard pass/fail metrics label it as passing despite practical inadequacy.

\subsection{L2: Interface Compatibility}

The generated code preserves the expected API surface---function signatures, struct field names, type definitions, and interface contracts.
An L2 failure occurs when semantically correct code cannot compile or integrate with the existing codebase due to surface-level mismatches.
This failure mode is particularly relevant for AI-generated code: LLMs optimize for semantic equivalence but do not prioritize interface consistency across independent generations.

A human developer, by contrast, naturally preserves interface conventions because they read and understand the codebase context. The agent operates without this contextual awareness, generating valid but incompatible interfaces in different runs.
A concrete example: an agent generating a REST API handler may produce a struct field named \texttt{UserID} in one run and \texttt{user\_id} in another. Both names are internally consistent, but if the test suite or framework expects a specific naming convention, the code fails at the serialization boundary despite being semantically identical. Such failures emerge from the misalignment between the agent's unbounded naming choices and the implicit interface contract assumed by the test suite or surrounding codebase.

\subsection{L3: Temporal Stability}

The solution produces consistent outcomes across repeated runs under identical conditions.
This layer captures non-deterministic behaviors: race conditions, timing sensitivity, scheduler-dependent outcomes, and error-handling gaps that only manifest probabilistically.
L3 failures are the most insidious because they are invisible to single-run evaluation---a task can pass once and fail repeatedly without any change to the code.
For example, a file-processing agent that opens a resource without deferred cleanup may succeed under sequential execution but fail when goroutine scheduling interleaves cleanup with subsequent operations. The generated code is identical across runs---only the runtime interleaving differs---yet the behavioral outcome varies in a way that single-run evaluation cannot capture.

\begin{figure}[h]
\centering
\begin{tabular}{p{0.45\textwidth}p{0.45\textwidth}}
\toprule
\textbf{Layer} & \textbf{Detectable by single run?} \\
\midrule
L1: Semantic correctness & Yes \\
L2: Interface compatibility & Yes (compilation) \\
L2: Interface compatibility & No (behavioral) \\
L3: Temporal stability & No \\
\bottomrule
\end{tabular}
\caption{Three failure layers and their detectability by single-run evaluation.}
\label{tab:layers}
\end{figure}

\section{Case Study: H2 Distributed Lock}

\subsection{Experimental Setup}

The H2 task required fixing a distributed lock deadlock bug in Go.
The correct implementation needed: (1) Lua-based atomic unlock with identity verification, (2) a watchdog goroutine for automatic TTL renewal, (3) thread-safe ownership management, and (4) proper error handling for Redis failures.

The agent was given the buggy implementation and a test suite with nine test cases covering mutex, identity verification, renewal, panic safety, concurrent locking (30 goroutines), and double-unlock safety.
The experiment was repeated five times from fresh agent sessions, each starting from the same buggy code.
All runs used a consistent temperature of $T = 0.7$ with default top-p and top-k sampling parameters, reflecting a realistic agent usage scenario where moderate output diversity is expected.
Each run was conducted in an isolated environment with a clean Go module cache and a fresh Docker-based Redis instance to prevent state contamination between trials.
The agent was prompted with an identical natural-language description of the deadlock bug and the expected fix strategy, without any in-run human intervention or guidance.
Run completion was capped at 50 tool invocations to bound experimental cost, and all five runs completed within this limit.
The full execution trace, including every file read, code generation step, compilation attempt, and test invocation, was logged for post-hoc failure analysis.

\subsection{Results}

The test suite evaluates the distributed lock implementation across nine scenarios: (T1) basic lock acquisition and release, (T2) identity verification during unlock, (T3) mutual exclusion under concurrent goroutines, (T4) automatic TTL renewal by the watchdog goroutine, (T5) renewal after extended hold beyond initial TTL, (T6) panic safety---recovering from a panic without state corruption, (T7) high-contention stress with 30 concurrent goroutines, (T8) double-unlock safety, and (T9) lock expiration with re-acquisition.
Tests T1--T6 and T8 are deterministic, while T7 and T9 depend on goroutine scheduling order and are inherently timing-sensitive.
This test suite structure was deliberately chosen to cover both functional correctness and concurrency robustness, providing diagnostic coverage across all three failure layers.

\begin{table}[h]
\centering
\caption{H2 results across five independent runs.}
\label{tab:h2}
\begin{tabular}{lclc}
\toprule
Run & Pass/Fail & Failure Layer & Tests Passed \\
\midrule
1 & Pass & --- & 6/6 \\
2 & Pass & --- & 6/6 \\
3 & Pass & --- & 6/6 \\
4 & Fail & L2 (compilation) & 0/6 \\
5 & Fail & L2+L3 (mixed) & 4/6 \\
\bottomrule
\end{tabular}
\end{table}

\subsection{Three-Layer Analysis}

\textbf{L1 --- Semantic Correctness (Intact):} All five implementations correctly implemented the distributed lock protocol---Lua-based atomic unlock, watchdog renewal, and ownership verification. The core locking semantics were correct in every run.

\textbf{L2 --- Interface Compatibility (Two Failures):} Runs 4 and 5 failed initially due to struct field naming mismatches. The test suite accessed \texttt{lock2.value} directly, but the agent used field names \texttt{id} (Run 4) and \texttt{ownerID} (Run 5) instead of the expected \texttt{value}. These naming differences caused compilation errors. While the behavior was semantically correct, the test suite's direct struct access was incompatible with the agent's generated API surface.

\textbf{L3 --- Temporal Stability (Partial Failure):} Run 5 additionally exhibited timing-sensitive behavior under concurrent load. The implementation discarded Lua script errors via blank identifier (\texttt{\_ = script.Run(...)}), causing transient Redis failures to be silently ignored. Under the concurrent stress test (30 goroutines), scheduling variability caused the watchdog renewal to occasionally miss its deadline, leading to premature lock expiry and partial test failure. This failure was non-deterministic---it depended on goroutine scheduling order.
The interaction between failure layers is noteworthy: Run 5's L2 naming mismatch, once corrected, revealed the underlying L3 instability that would otherwise have remained hidden beneath the compilation error. This masking effect suggests that post-hoc layer classification requires a two-phase diagnosis---first correcting purely mechanical interface issues, then assessing runtime stability---to avoid underestimating temporal instability in the presence of interface failures.

\section{Variance Index}

We propose a \textbf{Variance Index (VI)} for quantifying run-to-run reliability:

\[
VI = \frac{1}{n}\sum_{i=1}^{n}(x_i - \bar{x})^2
\]

where $x_i$ is the number of tests passed in the $i$-th run and $\bar{x}$ is the mean across $n$ runs.

\subsection{Empirical Values}

Across our variability analysis of six tasks with five runs each (30 total runs):

\begin{table}[h]
\centering
\caption{Variance Index across six tasks.}
\label{tab:vi}
\begin{tabular}{lcc}
\toprule
Task & Type & VI \\
\midrule
S1: Two-Sum & Sequential & 0.00 \\
S3: Binary search fix & Sequential & 0.00 \\
M1: LRU cache & Sequential & 0.00 \\
M4: SQL injection fix & Sequential & 0.00 \\
H1: Red-black tree & Sequential & 0.00 \\
H2: Distributed lock & Concurrent & \textbf{2.56} \\
\bottomrule
\end{tabular}
\end{table}

The five sequential tasks show $VI = 0.00$, confirming that zero variance is achievable for deterministic tasks with well-specified requirements.
The H2 concurrent task shows $VI = 2.56$, reflecting the probabilistic nature of its failures.

\subsection{Recommendation}

We recommend that AI coding agent evaluations report:
\begin{itemize}
\item Single-run pass rate (current standard)
\item Variance Index across 5--10 repeated runs
\item Layer-wise failure classification (L1/L2/L3) for failed runs
\end{itemize}

\section{Discussion}

\subsection{Implications for Evaluation Methodology}

Our findings suggest that the software engineering community should adopt multi-run evaluation as a standard practice for AI coding agents.
The Variance Index provides a simple, interpretable metric that complements pass rates.
For tasks with $VI > 0$, detailed failure analysis using the three-layer model can identify specific weaknesses.

\subsection{Implications for Test Suite Design}

L2 failures occurred because the test suite relied on direct struct field access rather than behavioral verification.
This suggests that test suites for AI evaluation should prefer interface-based verification methods---testing behavior through public APIs rather than internal structure---to reduce false negatives from naming inconsistencies.
Concretely, test harnesses should define API contracts through interfaces or abstract types against which the agent's output is compiled, decoupling structural naming choices from functional verification.
For concurrent tasks, evaluation suites should include load-variation stress tests that force scheduling variability, exposing L3 instabilities that are invisible under single-run conditions.
Multi-run test harnesses should further report flakiness ratios: the proportion of runs in which a given test transitions between pass and fail, as this metric distinguishes genuinely unstable behavior from isolated faults and provides a direct signal for temporal stability.
We also recommend a two-stage evaluation protocol---single-run screening followed by multi-run variability profiling for tasks that pass the initial screen---to balance evaluation cost against the need for reliability assessment.
Finally, test designers should consider whether their assertions test internal structure or external behavior; preferring the latter reduces the false positive rate from L2 naming variations while preserving coverage of genuine semantic errors.

\subsection{Limitations}

Our variability analysis covers six of eighteen tasks.
While the five sequential tasks showed zero variance, we cannot guarantee this generalizes to all deterministic tasks.
Similarly, while H2 showed high variance, other concurrent tasks may behave differently.
The five-run protocol provides a coarse variance estimate; tasks with low but nonzero variance would require 10--30 runs for reliable detection.
Future work should extend multi-run evaluation to larger, more diverse task sets and investigate whether the three-layer model generalizes to other agent architectures and programming languages.

\subsection{Toward Standardized Variability Reporting}

Based on our findings, we propose a minimal reporting standard for AI coding agent evaluations. Beyond the current single-run pass rate, evaluations should report: (1) the Variance Index across a minimum of five runs per task, (2) layer-wise failure classification for each failed run, and (3) the test suite's verification strategy (structural vs. behavioral) to contextualize L2 failure rates. Adopting this standard would enable meaningful cross-study comparisons of agent reliability and help the community identify task categories where current agents are most brittle.

\section{Conclusion}

We demonstrate that single-run pass rates can overestimate the reliability of AI coding agents.
Through a case study of a distributed lock task repeated five times, we identified three independent failure layers---semantic correctness, interface compatibility, and temporal stability---and observed that semantically correct code can fail at the interface or temporal layers in up to 40\% of runs.
The proposed three-layer model and variance index provide a more nuanced framework for evaluating AI-generated code.
We recommend that future benchmarks incorporate multi-run variability analysis as a standard practice for assessing AI coding agent reliability.

\begin{thebibliography}{00}
\bibitem{chen2021evaluating}
M. Chen \textit{et al.}, ``Evaluating large language models trained on code,'' \textit{arXiv:2107.03374}, 2021.
\bibitem{jimenez2023swebench}
C. E. Jimenez \textit{et al.}, ``SWE-bench: Can language models resolve real-world GitHub issues?'' \textit{arXiv:2310.06770}, 2023.
\bibitem{brown2020language}
T. Brown \textit{et al.}, ``Language models are few-shot learners,'' in \textit{NeurIPS}, 2020.
\bibitem{roziere2023code}
B. Roziere \textit{et al.}, ``Code Llama: Open foundation models for code,'' \textit{arXiv:2308.12950}, 2023.
\bibitem{li2023starcoder}
R. Li \textit{et al.}, ``StarCoder: may the source be with you!'' \textit{arXiv:2305.06161}, 2023.
\bibitem{austin2021program}
J. Austin \textit{et al.}, ``Program synthesis with large language models,'' \textit{arXiv:2108.07732}, 2021.
\bibitem{peng2023impact}
S. Peng \textit{et al.}, ``The impact of AI on developer productivity: Evidence from GitHub Copilot,'' \textit{arXiv:2302.06590}, 2023.
\bibitem{hendrycks2021measuring}
D. Hendrycks \textit{et al.}, ``Measuring coding challenge competence with APPS,'' \textit{NeurIPS}, 2021.
\bibitem{liu2023agentbench}
X. Liu \textit{et al.}, ``AgentBench: Evaluating LLMs as agents,'' in \textit{International Conference on Learning Representations (ICLR)}, 2024.
\end{thebibliography}

\end{document}
